\section{Step 58 Construct the full backward distance}\label{png-huff-s58}
Fill \texttt{huff\_distance\_extra}. Obtain the selected distance base and extra-bit count, retain the base at frame offset 40, request the ordinary numerical field and add it to the base. Keep the reconstructed distance in the temporary register used by the following bound checks.

Symbol zero has base one and no extra bits, so the distance is one. The zero-bit request must not fetch another byte. A later entry may add a nonzero field to its base; retaining only the symbol index would describe neither of those distances correctly.

\paragraph{Observe} Use \texttt{s58-1}; stop before \texttt{huff\_distance\_ready}. Inspect the state described above before continuing.

\Needspace{7\baselineskip}
\paragraph{Check} May a zero-extra-bit distance fetch a new input byte?
\begin{itemize}
  \item[A.] Yes, every distance consumes a byte.
  \item[B.] No; the zero-bit read retains the input cursor.
  \item[C.] Only after a final block.
\end{itemize}
